% Vietnamese translation for OI Public Library
% Source-File: IOI中国国家候选队论文/2005/杨弋/杨弋.ppt
\documentclass[11pt]{article}
\usepackage{oipl-vn}

\title{Bản trình chiếu: Ghi chú về quy hoạch động hình học}
\author{Yang Yi}
\date{2005}

\begin{document}
\maketitle

\origtitle{Slide companion for Yang Yi's geometric DP presentation}
\sourcepath{IOI China candidate team papers / 2005 / Yang Yi / Yang Yi.ppt}

\section{Phạm vi bản dịch}

Tệp gốc chỉ có PowerPoint. Phần chữ đọc được chứa các ký hiệu \(X_1,X_2\),
\(f(w,u,v)\), \(\operatorname{area}(x,k)\), bài \textit{Lock And Key}, giá trị
\(3000\), và độ phức tạp \(O(l^2mn^2)\), \(O(ln)\).

\section{Mô hình nhận diện được}

Các công thức \(f(w,u,v)\) và \(\operatorname{area}(x,k)\) cho thấy slide dùng
một hàm trạng thái phụ thuộc vào ba tham số hình học hoặc đoạn. Với bài
\textit{Lock And Key}, khả năng cao cần so khớp các miền hình học hoặc kiểm
tra giao nhau sau khi biến đổi.

\section{Tư tưởng thuật toán}

Khi bài hình học có nhiều tọa độ, ta thường giảm bài toán về các mốc biên,
diện tích, hoặc thứ tự góc. Nếu \(\operatorname{area}(x,k)\) tính được
\(O(1)\), bảng \(f\) có thể dùng để tránh tính lại các phần hình học lặp.

\section{Ghi chú}

Do phần văn bản chính không khôi phục đầy đủ, bản dịch này giữ vai trò
companion note, bảo toàn các ký hiệu và hướng thuật toán chắc chắn đọc được
từ slide.

\end{document}
